% !TeX root = ../sections/02_exercises.tex
\subsection{2-B-3.}
\begin{exercise}
	\[
	a_1=3,
	\qquad
	a_{n+1}
	=
	\frac{1}{2}
	\left(
	a_n+\frac{3}{a_n}
	\right)
	\]
	により定まる数列
	\[
	\{a_n\}_{n\in\mathbb{N}}
	\]
	の極限を求めよ。
\end{exercise}

\begin{background}
	この数列は，\(\sqrt{3}\) を求めるための反復法として現れる。
	漸化式
	\[
	a_{n+1}
	=
	\frac{1}{2}
	\left(
	a_n+\frac{3}{a_n}
	\right)
	\]
	は，正の数 \(a_n\) に対して，
	\(a_n\) と \(3/a_n\) の平均を取る形になっている。
	
	極限を求めるには，まず数列が収束することを示す必要がある。
	そのために，
	\[
	a_n>\sqrt{3}
	\]
	と
	\[
	a_{n+1}\leq a_n
	\]
	を示す。
	
	すなわち，\(\{a_n\}\) が下に有界な単調減少数列であることを示せば，
	単調収束定理より極限が存在する。
\end{background}

\begin{strategy}
	まず，すべての \(n\) に対して
	\[
	a_n>\sqrt{3}
	\]
	を示す。
	\(a_n>0\) のもとで計算すると，
	\[
	\begin{aligned}
		a_{n+1}-\sqrt{3}
		&=
		\frac{1}{2}
		\left(
		a_n+\frac{3}{a_n}
		\right)
		-\sqrt{3} \\
		&=
		\frac{a_n^2-2\sqrt{3}a_n+3}{2a_n} \\
		&=
		\frac{(a_n-\sqrt{3})^2}{2a_n}.
	\end{aligned}
	\]
	したがって，\(a_n>0\) ならば
	\[
	a_{n+1}\geq\sqrt{3}
	\]
	である。
	さらに初項 \(a_1=3>\sqrt{3}\) なので，
	帰納法により
	\[
	a_n>\sqrt{3}
	\]
	が従う。
	
	次に，単調減少性を示す。
	\[
	\begin{aligned}
		a_n-a_{n+1}
		&=
		a_n-
		\frac{1}{2}
		\left(
		a_n+\frac{3}{a_n}
		\right) \\
		&=
		\frac{a_n^2-3}{2a_n}.
	\end{aligned}
	\]
	すでに \(a_n>\sqrt{3}\) が分かっているので，
	\[
	a_n^2-3>0
	\]
	である。
	したがって
	\[
	a_n-a_{n+1}>0
	\]
	となり，
	\[
	a_{n+1}<a_n
	\]
	である。
	
	よって，\(\{a_n\}\) は下に有界な単調減少数列である。
	したがって収束する。
	その極限を
	\[
	\alpha=\lim_{n\to\infty}a_n
	\]
	とおく。
	
	漸化式の両辺で \(n\to\infty\) とすると，
	\[
	\alpha
	=
	\frac{1}{2}
	\left(
	\alpha+\frac{3}{\alpha}
	\right)
	\]
	である。
	これを整理すると
	\[
	2\alpha=\alpha+\frac{3}{\alpha}
	\]
	より
	\[
	\alpha=\frac{3}{\alpha}.
	\]
	したがって
	\[
	\alpha^2=3.
	\]
	また \(a_n>\sqrt{3}>0\) なので \(\alpha\geq0\) である。
	よって
	\[
	\alpha=\sqrt{3}.
	\]
\end{strategy}